\documentclass[11pt]{article}

\usepackage[margin=1in]{geometry}
\usepackage{amsmath, amssymb, amsthm}
\usepackage{booktabs}
\usepackage{hyperref}
\usepackage{xcolor}
\usepackage{enumitem}
\usepackage{longtable}
\usepackage{array}

\hypersetup{
  colorlinks=true,
  linkcolor=blue,
  urlcolor=blue,
  citecolor=blue
}

\title{Formalizing Finite-Class Learning Guarantees in Lean 4\\
with an MA-LoT-Inspired Multi-LLM Proof Repair Workflow}
\author{Dhruv}
\date{}

\newcommand{\Hclass}{\mathcal{H}}
\newcommand{\Ddist}{\mathcal{D}}
\newcommand{\Xspace}{\mathcal{X}}
\newcommand{\Yspace}{\mathcal{Y}}
\newcommand{\LD}{L_{\Ddist}}
\newcommand{\LS}{L_S}
\newcommand{\hh}{\widehat{h}}
\newcommand{\hstar}{h^\star}
\newcommand{\eps}{\epsilon}
\newcommand{\ind}{\mathbf{1}}

\theoremstyle{definition}
\newtheorem{definition}{Definition}
\newtheorem{theorem}{Theorem}
\newtheorem{lemma}{Lemma}
\newtheorem{remark}{Remark}

\begin{document}

\maketitle

\begin{abstract}
This project formalizes finite hypothesis class learning guarantees in Lean 4
and studies an MA-LoT-inspired workflow for repairing AI-generated Lean proofs.
The mathematical topic is empirical risk minimization for finite hypothesis
classes in the realizable and agnostic settings. The strongest fully verified
Lean result is the deterministic agnostic ERM theorem: under uniform
convergence, ERM competes with a comparator hypothesis up to an additive
$2\eps$ term. The realizable and high-probability agnostic guarantees are
represented using explicit assumed concentration results rather than by
formalizing probability theory from scratch.

The workflow contribution documents how Codex and Claude generate proof
attempts from a shared theorem prompt, how Lean verifier feedback exposes
errors, how repair iterations fix those errors, and how comparison across
models produces both human-facing explanations and model-facing repair
instructions.
\end{abstract}

\tableofcontents

\section{Reader's Guide}

This report is organized around the three requirements in the final project
instructions.

\begin{enumerate}[leftmargin=*]
  \item \textbf{Formalized proof.} The Lean code is in
  \texttt{LearningTheoryProject/}. The main fully proved theorem is
  \texttt{agnostic\_erm\_deterministic} in
  \texttt{LearningTheoryProject/Agnostic.lean}. Probability concentration
  results are stated explicitly as assumptions in
  \texttt{LearningTheoryProject/AssumedResults.lean}.
  \item \textbf{Designed proof presentation.} The mathematical proof is
  explained in this report, especially in the realizable and agnostic proof
  sections. Additional reader-facing materials are in \texttt{notes/} and
  \texttt{presentation/}, including an annotated proof walkthrough, dependency
  graph, paper-versus-Lean table, and final slide deck.
  \item \textbf{AI workflow.} The MA-LoT-inspired workflow is documented in
  \texttt{ai\_workflow/}, \texttt{malot\_workflow/},
  \texttt{multi\_model\_workflow/}, \texttt{generated\_proofs/}, and
  \texttt{proof\_refactoring/}. These files record prompts, generated proof
  attempts, Lean verifier feedback, repair iterations, comparison reports, and
  reflection.
\end{enumerate}

The most important formalization boundary is this: the deterministic agnostic
ERM theorem is fully proved in Lean, while probability concentration results
such as Hoeffding-style uniform convergence are stated as explicit assumptions.

\section{Main Result Summary}

\begin{center}
\begin{tabular}{p{0.27\textwidth}p{0.34\textwidth}p{0.29\textwidth}}
\toprule
\textbf{Result} & \textbf{Mathematical meaning} & \textbf{Lean status} \\
\midrule
Realizable guarantee
& Consistent hypotheses have small true risk with high probability.
& Wrapper theorem using an assumed concentration result. \\
\addlinespace
Agnostic deterministic theorem
& Under uniform convergence, ERM is within $2\eps$ of a comparator.
& Fully proved in Lean. \\
\addlinespace
Agnostic high-probability guarantee
& ERM is close to the best hypothesis in the class with high probability.
& Uses assumed uniform convergence plus the verified deterministic theorem. \\
\bottomrule
\end{tabular}
\end{center}

\section{Assumptions and Scope at a Glance}

The project deliberately separates mathematical assumptions, Lean assumptions,
and implementation choices.

\subsection{Mathematical Assumptions}

\begin{itemize}[leftmargin=*]
  \item The hypothesis class $\Hclass$ is finite.
  \item Training examples are drawn independently from an unknown distribution
  $\Ddist$.
  \item In the realizable setting, there exists $\hstar\in\Hclass$ such that
  $\LD(\hstar)=0$.
  \item In the agnostic setting, no perfect classifier is assumed. The learner is
  compared with a best-in-class hypothesis $\hstar$.
  \item The agnostic deterministic theorem assumes uniform convergence: empirical
  risk and true risk are within $\eps$ for every hypothesis in the class.
\end{itemize}

\subsection{Assumptions Represented Explicitly in Lean}

The file \texttt{LearningTheoryProject/AssumedResults.lean} contains the
probability-facing assumptions. The key assumed interfaces are:

\begin{itemize}[leftmargin=*]
  \item \texttt{HighProbabilityEvent}: an abstract marker for ``with probability
  at least $1-\delta$.''
  \item \texttt{highProbability\_mono}: an assumed monotonicity principle for
  high-probability events.
  \item \texttt{finiteClass\_realizable\_concentration}: the assumed
  finite-class concentration result for the realizable setting.
  \item \texttt{finiteClass\_agnostic\_uniform\_convergence}: the assumed
  finite-class uniform convergence result for the agnostic setting.
\end{itemize}

\subsection{Implementation Scope Choices}

\begin{itemize}[leftmargin=*]
  \item Risks are represented as abstract \texttt{Nat}-valued quantities in Lean.
  This keeps the deterministic order proof lightweight.
  \item The usual logarithmic sample complexity rates are proved and explained
  mathematically in this report, but the Lean code represents the sample-size
  requirements through abstract sample-condition predicates.
  \item The project does not formalize measure theory, probability spaces,
  Hoeffding's inequality, or the finite-class union bound from first principles.
  \item The project is MA-LoT-inspired. It does not implement the full MA-LoT
  system.
\end{itemize}

\section{Project Overview}

The project has two connected goals.

\begin{enumerate}[leftmargin=*]
  \item \textbf{Lean formalization goal.} Formalize the structure of finite
  hypothesis class learning guarantees in Lean 4. The project covers the
  realizable setting and the agnostic setting, while keeping probability-heavy
  concentration facts explicit as assumptions.
  \item \textbf{Proof workflow goal.} Study an MA-LoT-inspired proof repair
  workflow in which AI-generated Lean proofs are checked by Lean, repaired from
  verifier feedback, compared across models, simplified only after verification,
  and translated into human-readable explanations.
\end{enumerate}

The central communication theme is:

\begin{quote}
Verification is not the same as understanding.
\end{quote}

Lean can verify that a proof term is correct. It does not automatically explain
why the theorem matters, what informal assumptions were hidden in the paper
proof, or why an AI-generated proof failed. The repository therefore includes
both machine-checkable Lean code and human-readable proof artifacts.

\section{Motivation: Why Assumptions Matter}

No Free Lunch-style results say that learning is impossible without assumptions.
If every possible labeling function is allowed, then no learning algorithm is
universally better than another. This does not contradict finite-class learning
guarantees. Instead, it explains why assumptions are necessary.

This project studies a classical positive result after adding a concrete
assumption:

\begin{quote}
The hypothesis class $\Hclass$ is finite.
\end{quote}

Finiteness makes learning possible in the standard proof because it allows a
union bound over all hypotheses. In the realizable case, the proof rules out bad
hypotheses that nevertheless fit the sample. In the agnostic case, the proof
uses uniform convergence to relate empirical risk to true risk for every
hypothesis in the class.

\section{Mathematical Setup}

\subsection{Binary Classification}

Let $\Xspace$ be an input space and let $\Yspace=\{0,1\}$ be the label space. A
hypothesis is a classifier
\[
  h : \Xspace \to \Yspace.
\]

The hypothesis class is finite:
\[
  \Hclass = \{h_1,h_2,\ldots,h_N\}.
\]

A data point is
\[
  z=(x,y),
\]
where $x\in\Xspace$ and $y\in\Yspace$. The data distribution is an unknown
distribution $\Ddist$ over $\Xspace\times\Yspace$.

A training sample is
\[
  S=(z_1,z_2,\ldots,z_m),
\]
where each $z_i=(x_i,y_i)$ is drawn independently from $\Ddist$.

\subsection{True Risk and Empirical Risk}

\begin{definition}[True risk]
For a hypothesis $h$, the true risk is
\[
  \LD(h) = \Pr_{(x,y)\sim\Ddist}[h(x)\ne y].
\]
This is the error probability on a fresh random example from the data
distribution.
\end{definition}

\begin{definition}[Empirical risk]
For a hypothesis $h$ and sample $S$, the empirical risk is
\[
  \LS(h) =
  \frac{1}{m}\sum_{i=1}^{m}\ind[h(x_i)\ne y_i].
\]
This is the fraction of training examples that $h$ classifies incorrectly.
\end{definition}

\begin{definition}[Empirical risk minimization]
An empirical risk minimizer is a hypothesis
\[
  \hh \in \arg\min_{h\in\Hclass}\LS(h).
\]
Equivalently,
\[
  \forall h\in\Hclass,\quad \LS(\hh)\le \LS(h).
\]
\end{definition}

\section{Realizable Setting}

\subsection{Assumption}

The realizable setting assumes that the hypothesis class contains a perfect
classifier:
\[
  \exists \hstar\in\Hclass
  \quad\text{such that}\quad
  \LD(\hstar)=0.
\]

In words, at least one hypothesis in the class makes no mistakes under the true
data distribution.

\subsection{Consistency of ERM}

If $\LD(\hstar)=0$, then with probability one $\hstar$ makes no mistakes on the
training sample. Therefore
\[
  \LS(\hstar)=0.
\]

Since $\hh$ is an empirical risk minimizer,
\[
  \LS(\hh)\le \LS(\hstar)=0.
\]

Empirical risk is nonnegative, so
\[
  \LS(\hh)=0.
\]

Thus, in the realizable setting, ERM returns a hypothesis that is consistent
with the sample.

\subsection{Bad Hypotheses}

Fix an accuracy parameter $\eps>0$. Define the set of bad hypotheses:
\[
  B = \{h\in\Hclass : \LD(h)>\eps\}.
\]

If ERM fails to return a good hypothesis, then $\LD(\hh)>\eps$. But in the
realizable setting ERM is consistent, so this means some bad hypothesis must
have empirical risk zero. Therefore
\[
  \{\LD(\hh)>\eps\}
  \subseteq
  \{\exists h\in B : \LS(h)=0\}.
\]

Taking probabilities gives
\[
  \Pr[\LD(\hh)>\eps]
  \le
  \Pr[\exists h\in B : \LS(h)=0].
\]

\subsection{One Bad Hypothesis}

Fix $h\in B$. Since $h$ is bad,
\[
  \LD(h)>\eps.
\]

This means the probability that $h$ makes a mistake on one fresh sample is more
than $\eps$. Therefore the probability that $h$ is correct on one sample is less
than $1-\eps$.

For $h$ to have zero empirical risk, it must be correct on all $m$ independent
examples. Hence
\[
  \Pr[\LS(h)=0]
  =
  (1-\LD(h))^m
  \le
  (1-\eps)^m.
\]

Using the standard inequality $1-\eps\le e^{-\eps}$, we obtain
\[
  \Pr[\LS(h)=0]\le e^{-\eps m}.
\]

\subsection{Union Bound}

Using the union bound over all bad hypotheses,
\[
\begin{aligned}
  \Pr[\exists h\in B : \LS(h)=0]
  &\le \sum_{h\in B}\Pr[\LS(h)=0] \\
  &\le |B|e^{-\eps m} \\
  &\le |\Hclass|e^{-\eps m}.
\end{aligned}
\]

Thus
\[
  \Pr[\LD(\hh)>\eps]\le |\Hclass|e^{-\eps m}.
\]

\subsection{Sample Complexity}

To make the failure probability at most $\delta$, it is enough to require
\[
  |\Hclass|e^{-\eps m}\le \delta.
\]

Taking logarithms:
\[
  \log|\Hclass|-\eps m\le \log\delta.
\]

Rearranging:
\[
  \eps m \ge \log|\Hclass|+\log\frac{1}{\delta}.
\]

Therefore it is enough that
\[
  m\ge
  \frac{1}{\eps}
  \left(
    \log|\Hclass|+\log\frac{1}{\delta}
  \right).
\]

\begin{theorem}[Realizable finite-class guarantee, informal]
If $\Hclass$ is finite, the setting is realizable, and
\[
  m\ge
  \frac{1}{\eps}
  \left(
    \log|\Hclass|+\log\frac{1}{\delta}
  \right),
\]
then with probability at least $1-\delta$, every hypothesis in $\Hclass$ that is
consistent with the sample has true risk at most $\eps$.
\end{theorem}

\subsection{Lean Status}

In Lean, the probability-heavy realizable concentration theorem is stated as an
explicit assumption:
\[
  \texttt{finiteClass\_realizable\_concentration}.
\]

The theorem wrapper is:
\[
  \texttt{realizable\_learning\_guarantee}.
\]

It is located in:
\[
  \texttt{LearningTheoryProject/Realizable.lean}.
\]

The Lean theorem does not reprove Hoeffding, the union bound, or the logarithmic
sample complexity algebra. Instead, it formalizes the surrounding theorem shape
using a named concentration assumption.

\section{Agnostic Setting}

\subsection{Assumption Removed}

The agnostic setting removes the assumption that a perfect hypothesis exists.
There may be no $h\in\Hclass$ with $\LD(h)=0$.

Instead, define $\hstar$ to be a best hypothesis in $\Hclass$:
\[
  \hstar \in \arg\min_{h\in\Hclass}\LD(h).
\]

The goal is not perfection. The goal is to show that ERM performs almost as well
as $\hstar$.

\subsection{Uniform Convergence}

The key deterministic assumption is uniform convergence:
\[
  \forall h\in\Hclass,\quad |\LD(h)-\LS(h)|\le \eps.
\]

This absolute-value statement gives two directed inequalities:
\[
  \LD(h)\le \LS(h)+\eps,
\]
and
\[
  \LS(h)\le \LD(h)+\eps.
\]

Lean represents this directly as two inequalities rather than as an absolute
value. This is one example of Lean forcing a paper proof to become more
explicit.

\subsection{Deterministic Agnostic ERM Lemma}

\begin{lemma}[Deterministic agnostic ERM guarantee]
Assume uniform convergence and assume $\hh$ is an empirical risk minimizer. Then
for any comparator $\hstar$,
\[
  \LD(\hh)\le \LD(\hstar)+2\eps.
\]
\end{lemma}

\begin{proof}
Start with the true risk of $\hh$. By uniform convergence applied to $\hh$,
\[
  \LD(\hh)\le \LS(\hh)+\eps.
\]

Since $\hh$ is an empirical risk minimizer,
\[
  \LS(\hh)\le \LS(\hstar).
\]

Adding $\eps$ to both sides gives
\[
  \LS(\hh)+\eps\le \LS(\hstar)+\eps.
\]

By uniform convergence applied to $\hstar$ in the other direction,
\[
  \LS(\hstar)\le \LD(\hstar)+\eps.
\]

Adding $\eps$ gives
\[
  \LS(\hstar)+\eps
  \le
  \LD(\hstar)+\eps+\eps.
\]

Combining the inequalities:
\[
\begin{aligned}
  \LD(\hh)
  &\le \LS(\hh)+\eps \\
  &\le \LS(\hstar)+\eps \\
  &\le \LD(\hstar)+\eps+\eps \\
  &= \LD(\hstar)+2\eps.
\end{aligned}
\]

Thus ERM is within $2\eps$ of the comparator hypothesis.
\end{proof}

The two $\eps$ terms come from applying uniform convergence twice: once for
$\hh$ and once for $\hstar$.

\subsection{Lean Status}

This deterministic result is the strongest fully verified theorem in the
project. In Lean it is named:
\[
  \texttt{agnostic\_erm\_deterministic}.
\]

It is located in:
\[
  \texttt{LearningTheoryProject/Agnostic.lean}.
\]

The Lean proof follows the same chain:
\begin{itemize}[leftmargin=*]
  \item apply uniform convergence to \texttt{hHat},
  \item use the ERM inequality comparing \texttt{hHat} to \texttt{hStar},
  \item add \texttt{eps} to both sides with \texttt{Nat.add\_le\_add\_right},
  \item apply uniform convergence to \texttt{hStar},
  \item rewrite associativity with \texttt{Nat.add\_assoc}.
\end{itemize}

\subsection{High-Probability Agnostic Bound}

For a fixed hypothesis $h$, empirical risk is an average of independent
zero-one random variables:
\[
  \ind[h(x_i)\ne y_i].
\]

Its expectation is the true risk:
\[
  \mathbb{E}[\LS(h)] = \LD(h).
\]

Hoeffding's inequality gives
\[
  \Pr\left[|\LS(h)-\LD(h)|>\eps\right]
  \le
  2e^{-2m\eps^2}.
\]

Using the union bound over $\Hclass$:
\[
\begin{aligned}
  \Pr\left[
    \exists h\in\Hclass :
    |\LS(h)-\LD(h)|>\eps
  \right]
  &\le
  \sum_{h\in\Hclass}
  \Pr\left[|\LS(h)-\LD(h)|>\eps\right] \\
  &\le
  2|\Hclass|e^{-2m\eps^2}.
\end{aligned}
\]

To make this failure probability at most $\delta$, require
\[
  2|\Hclass|e^{-2m\eps^2}\le \delta.
\]

Solving gives
\[
  m\ge
  \frac{1}{2\eps^2}
  \log\frac{2|\Hclass|}{\delta}.
\]

Therefore, with probability at least $1-\delta$,
\[
  \LD(\hh)
  \le
  \min_{h\in\Hclass}\LD(h)+2\eps.
\]

In Lean, the finite-class uniform convergence theorem is assumed as:
\[
  \texttt{finiteClass\_agnostic\_uniform\_convergence}.
\]

The high-probability wrapper theorem is:
\[
  \texttt{agnostic\_high\_probability\_guarantee}.
\]

\section{Why the Rates Differ}

The realizable rate is:
\[
  m =
  O\left(
    \frac{\log|\Hclass|+\log(1/\delta)}{\eps}
  \right).
\]

The agnostic rate is:
\[
  m =
  O\left(
    \frac{\log|\Hclass|+\log(1/\delta)}{\eps^2}
  \right).
\]

The difference comes from the type of event being bounded.

In the realizable setting, the proof asks whether a bad hypothesis makes zero
mistakes on all samples. That probability decays like
\[
  e^{-m\eps}.
\]

In the agnostic setting, the proof asks whether empirical averages are close to
expectations for all hypotheses. Hoeffding's inequality gives decay like
\[
  e^{-m\eps^2}.
\]

Solving these exponential bounds for $m$ gives the different rates.

\section{Lean Formalization}

\subsection{Design Choice: Abstract Risk}

The Lean project uses:
\[
  \texttt{abbrev Risk := Nat}.
\]

This is an abstraction. It does not claim that risks are literally natural
numbers. The purpose is to focus on the deterministic order reasoning needed for
the ERM theorem while avoiding a full formalization of real-valued probability.

\subsection{Main Definitions}

The core structure is:
\[
  \texttt{LearningProblem Hypothesis}.
\]

It contains:
\begin{itemize}[leftmargin=*]
  \item \texttt{trueRisk : Hypothesis -> Risk},
  \item \texttt{empiricalRisk : Hypothesis -> Risk},
  \item \texttt{classMember : Hypothesis -> Prop},
  \item \texttt{classSize : Nat}.
\end{itemize}

The main predicates are:
\begin{itemize}[leftmargin=*]
  \item \texttt{InClass},
  \item \texttt{NonemptyClass},
  \item \texttt{Consistent},
  \item \texttt{BadHypothesis},
  \item \texttt{Realizable},
  \item \texttt{UniformConvergence},
  \item \texttt{EmpiricalRiskMinimizer},
  \item \texttt{TrueRiskOptimal}.
\end{itemize}

\subsection{Lean Theorem Interfaces and Assumptions}

The Lean project distinguishes ordinary definitions, proved deterministic
lemmas, and assumed probability interfaces.

\paragraph{Definitions.}
The following objects are definitions, not assumptions:
\begin{itemize}[leftmargin=*]
  \item \texttt{Consistent P h}: the empirical risk of \texttt{h} is zero.
  \item \texttt{BadHypothesis P eps h}: the true risk of \texttt{h} is greater
  than \texttt{eps}.
  \item \texttt{Realizable P}: there exists a class member with zero true risk.
  \item \texttt{UniformConvergence P eps}: both directed risk inequalities hold
  for every hypothesis.
  \item \texttt{EmpiricalRiskMinimizer P hHat}: \texttt{hHat} has empirical risk
  no larger than any comparator.
  \item \texttt{TrueRiskOptimal P hStar}: \texttt{hStar} has true risk no larger
  than any comparator.
\end{itemize}

\paragraph{Proved deterministic helper lemmas.}
The following are proved directly from definitions:
\begin{itemize}[leftmargin=*]
  \item \texttt{uniform\_true\_le\_empirical\_plus},
  \item \texttt{uniform\_empirical\_le\_true\_plus},
  \item \texttt{consistent\_empirical\_zero},
  \item \texttt{class\_member\_iff\_inClass}.
\end{itemize}

\paragraph{Assumed probability interfaces.}
The following are axioms in \texttt{AssumedResults.lean}:
\begin{itemize}[leftmargin=*]
  \item \texttt{HighProbabilityEvent}: represents an event holding with
  probability at least $1-\delta$.
  \item \texttt{highProbability\_mono}: if event $p$ implies event $q$, then a
  high-probability guarantee for $p$ gives a high-probability guarantee for $q$.
  \item \texttt{finiteClass\_realizable\_concentration}: the assumed finite-class
  realizable concentration theorem.
  \item \texttt{finiteClass\_agnostic\_uniform\_convergence}: the assumed
  finite-class agnostic uniform convergence theorem.
\end{itemize}

\paragraph{Main proved theorem.}
The theorem \texttt{agnostic\_erm\_deterministic} is fully proved in Lean. It
uses no probability axiom directly; it uses only uniform convergence, ERM, and
elementary natural-number order reasoning.

\paragraph{Wrapper theorems.}
The theorems \texttt{realizable\_learning\_guarantee} and
\texttt{agnostic\_high\_probability\_guarantee} compile in Lean but depend on
the assumed concentration interfaces. They are included to show the final
learning-theory theorem shape while keeping the probability boundary explicit.

\subsection{Formalization Boundary}

The project separates proved results from assumed results.

\begin{longtable}{>{\raggedright\arraybackslash}p{0.42\textwidth}
                  >{\raggedright\arraybackslash}p{0.48\textwidth}}
\toprule
\textbf{Component} & \textbf{Status} \\
\midrule
\endhead
Basic learning definitions & Formalized in Lean. \\
Consistency, realizability, ERM, uniform convergence & Formalized in Lean. \\
Deterministic agnostic ERM theorem & Fully proved in Lean. \\
Agnostic high-probability wrapper & Proved using an assumed concentration theorem. \\
Realizable learning guarantee wrapper & Proved using an assumed concentration theorem. \\
Hoeffding inequality & Assumed/scoped out. \\
Finite-class probability union bound & Assumed/scoped out. \\
Exact logarithmic sample complexity in Lean & Explained mathematically, not fully formalized. \\
Full MA-LoT system & Not implemented; workflow is MA-LoT-inspired. \\
\bottomrule
\end{longtable}

\section{Repository Layout}

The repository is organized so that the Lean code, mathematical notes,
presentation materials, AI workflow artifacts, and final report are separated.

\subsection{Lean Files}

\begin{longtable}{>{\raggedright\arraybackslash}p{0.36\textwidth}
                  >{\raggedright\arraybackslash}p{0.54\textwidth}}
\toprule
\normalfont\textbf{File} & \textbf{Purpose} \\
\midrule
\endhead
\small LearningTheoryProject/BasicDefinitions.lean
& Defines the abstract learning vocabulary: risk, classifiers, learning problems,
class membership, consistency, realizability, ERM, and uniform convergence. \\
\small LearningTheoryProject/AssumedResults.lean
& States the high-probability interface and assumed concentration results. \\
\small LearningTheoryProject/Realizable.lean
& Contains helper lemmas and the realizable theorem wrapper. \\
\small LearningTheoryProject/Agnostic.lean
& Contains the fully proved deterministic agnostic ERM theorem and the
high-probability agnostic wrapper. \\
\small LearningTheoryProject/DemoExamples.lean
& Contains small examples used in the live presentation. \\
\small LearningTheoryProject/MainTheorem.lean
& Imports the complete formalization and describes the main results. \\
\bottomrule
\end{longtable}

\subsection{Notes and Proof Explanation}

\begin{itemize}[leftmargin=*]
  \item \texttt{notes/math\_overview.md}: overview of the learning problem.
  \item \texttt{notes/proof\_realizable.md}: realizable proof explanation.
  \item \texttt{notes/proof\_agnostic.md}: agnostic proof explanation.
  \item \texttt{notes/lean\_translation\_notes.md}: paper proof to Lean translation notes.
\end{itemize}

\subsection{Presentation Materials}

\begin{itemize}[leftmargin=*]
  \item \texttt{presentation/final\_slides.pptx}: final delivered slide deck.
  \item \texttt{presentation/speaker\_notes.md}: slide-by-slide speaking script.
  \item \texttt{presentation/annotated\_proof\_walkthrough.md}: Lean proof next to human explanation.
  \item \texttt{presentation/paper\_vs\_lean\_table.md}: informal proof phrases versus explicit Lean requirements.
  \item \texttt{presentation/proof\_dependency\_graph.md}: dependency diagram for definitions, assumptions, and theorems.
  \item \texttt{presentation/proof\_pipeline\_diagram.md}: proof generation and repair workflow.
  \item \texttt{presentation/visual\_theorem\_flow.md}: theorem-flow explanation.
\end{itemize}

\section{MA-LoT-Inspired Proof Repair Workflow}

\subsection{What MA-LoT Contributes Conceptually}

MA-LoT motivates a theorem-proving workflow where natural-language reasoning,
Lean verification, and proof repair interact. The important lesson for this
project is that an LLM-generated proof is not trusted directly. It is checked by
Lean, and Lean's feedback is used to guide correction.

This project is explicitly MA-LoT-inspired. It does not implement the full
MA-LoT system.

\subsection{Adapted Workflow}

The adapted workflow is:

\begin{center}
\begin{tabular}{c}
English theorem prompt \\
$\downarrow$ \\
Codex proof attempt and Claude proof attempt \\
$\downarrow$ \\
Lean verifier feedback \\
$\downarrow$ \\
repair iteration \\
$\downarrow$ \\
comparison and reconciliation \\
$\downarrow$ \\
verified Lean theorem \\
$\downarrow$ \\
human-readable explanation
\end{tabular}
\end{center}

\subsection{Why This Matters}

LLMs can generate plausible-looking proofs that are not correct Lean proofs.
Lean catches errors such as:
\begin{itemize}[leftmargin=*]
  \item missing assumptions,
  \item wrong inequality directions,
  \item hallucinated lemma names,
  \item arithmetic gaps,
  \item tactic failures,
  \item theorem statements that are too informal.
\end{itemize}

The workflow turns those errors into repair data.

\subsection{AI Collaboration Reflection}

Across the project, AI assistance was most useful for proposing proof structure,
drafting explanatory text, and producing alternative Lean attempts that could be
checked against the verifier. The collaboration evolved in three stages.

\paragraph{Planning stage.}
AI assistance helped turn the broad project idea into a scoped plan: formalize
finite-class learning, prioritize the deterministic agnostic theorem, isolate
probability facts as assumptions, and build explanation artifacts in parallel
with Lean code.

\paragraph{Proof stage.}
AI-generated proof attempts were useful as drafts, but not as final authority.
The most important rule was that no proof was treated as correct until Lean
accepted it. The AI was delegated proof-attempt generation, proof repair
suggestions, and explanation drafting. Lean was the verifier.

\paragraph{Presentation stage.}
AI assistance helped convert formal Lean statements into human-readable proof
walkthroughs, diagrams, speaker notes, and the final report. This was especially
useful for explaining what Lean forced to become explicit: directed inequalities,
named assumptions, and arithmetic rewrites.

The workflow also showed where AI can mislead. It sometimes used theorem names
that did not exist, treated a conjunction as if it were a single inequality,
tried automation that could not see the right constraints, or skipped arithmetic
steps that a paper proof would call ``obvious.'' These failures became useful
only after Lean exposed them and they were logged, diagnosed, and repaired.

\section{Multi-Model Comparison Workflow}

\subsection{Blind Generation}

Codex and Claude are given the same theorem prompt:
\[
  \texttt{theorem\_prompts/agnostic\_erm\_theorem.md}.
\]

The final verified Lean proof is not given to them during blind generation. This
prevents the models from merely copying the final proof and makes the comparison
more meaningful.

\subsection{Generated Proof Artifacts}

The generated files are stored in:
\begin{itemize}[leftmargin=*]
  \item \texttt{generated\_proofs/codex/agnostic\_attempt\_1.lean},
  \item \texttt{generated\_proofs/codex/agnostic\_attempt\_2\_repaired.lean},
  \item \texttt{generated\_proofs/codex/simplified\_attempt.lean},
  \item \texttt{generated\_proofs/claude/agnostic\_attempt\_1.lean},
  \item \texttt{generated\_proofs/claude/agnostic\_attempt\_2\_repaired.lean},
  \item \texttt{generated\_proofs/claude/simplified\_attempt.lean}.
\end{itemize}

\subsection{Comparison Script}

The script
\[
  \texttt{multi\_model\_workflow/run\_comparison.py}
\]
checks each generated file with Lean and records structural features.

It compares:
\begin{itemize}[leftmargin=*]
  \item whether the file compiles,
  \item Lean's exit code,
  \item verifier error output,
  \item line count,
  \item SHA hash prefix,
  \item whether the proof uses \texttt{calc},
  \item whether the proof uses \texttt{omega},
  \item whether it uses \texttt{Nat.add\_le\_add\_right},
  \item whether it uses \texttt{Nat.add\_assoc},
  \item whether it calls the main theorem directly,
  \item whether it extracts conjunction directions using \texttt{.left} or \texttt{.right}.
\end{itemize}

The comparison is not a deep semantic equivalence checker. Lean verifies whether
each proof checks against its theorem statement. The script then compares proof
structure and usefulness for explanation or simplification.

\subsection{Repair and Simplification}

The comparison script generates:
\begin{itemize}[leftmargin=*]
  \item \texttt{multi\_model\_workflow/comparison\_report.md},
  \item \texttt{multi\_model\_workflow/repair\_recommendations\_for\_models.md}.
\end{itemize}

The repair recommendations tell Codex and Claude what to borrow from each other.
For example:
\begin{itemize}[leftmargin=*]
  \item Codex should explicitly extract uniform convergence directions.
  \item Claude should avoid relying on \texttt{omega} for abstract
  \texttt{Nat}-valued risk functions.
  \item Both should preserve a readable \texttt{calc} chain for explanation.
\end{itemize}

Simplification is only accepted after Lean verification. A shorter proof is not
kept merely because it looks cleaner; it must still pass \texttt{lake build}.

\subsection{How Comparison Prevents Shortcuts}

The comparison workflow prevents shortcut behavior in three ways.

\begin{enumerate}[leftmargin=*]
  \item Both models start from the same theorem prompt, not from the final proof.
  \item Lean checks whether each generated proof actually compiles.
  \item The comparison exposes missing assumptions, hallucinated lemmas, tactic
  failures, and proofs that only sound correct in English.
\end{enumerate}

Thus, the LLM is a proof assistant, but Lean remains the verifier.

\section{Proof Repair Case Study}

The agnostic ERM theorem was the main repair case study.

\subsection{Failure Pattern 1: Uniform Convergence Treated Too Informally}

Informally, the proof says ``by uniform convergence.'' In Lean, uniform
convergence is a conjunction with two directed inequalities. A failed proof
attempt tried to use the whole uniform convergence hypothesis directly.

The repair was to extract the needed direction:
\[
  \texttt{(hUC hHat).left}
\]
or
\[
  \texttt{(hUC hStar).right}.
\]

\subsection{Failure Pattern 2: ERM Used Before Matching the Goal}

ERM gives:
\[
  \LS(\hh)\le \LS(\hstar).
\]

But the proof goal needs:
\[
  \LS(\hh)+\eps\le \LS(\hstar)+\eps.
\]

The repair was:
\[
  \texttt{Nat.add\_le\_add\_right hERMStar eps}.
\]

\subsection{Failure Pattern 3: Arithmetic Hidden by ``Rearranging''}

The paper proof says the result follows by rearranging. Lean requires the exact
final expression:
\[
  \texttt{P.trueRisk hStar + (eps + eps)}.
\]

The repair used:
\[
  \texttt{rw [Nat.add\_assoc]}.
\]

\section{Verification and Reproducibility}

The project was verified with Lean 4.29.1 and Lake 5.0.0.

\subsection{Build Command}

To verify the Lean project:

\begin{verbatim}
lake build
\end{verbatim}

Expected output:

\begin{verbatim}
Build completed successfully (16 jobs).
\end{verbatim}

\subsection{Comparison Command}

To reproduce the multi-model comparison:

\begin{verbatim}
python multi_model_workflow/run_comparison.py
\end{verbatim}

Expected output:

\begin{verbatim}
Wrote multi_model_workflow\comparison_report.md
Wrote multi_model_workflow\repair_recommendations_for_models.md
\end{verbatim}

\section{Limitations}

\begin{itemize}[leftmargin=*]
  \item Risks are abstract \texttt{Nat} values, not real-valued probabilities.
  \item Hoeffding's inequality is not formalized in Lean.
  \item The finite-class probability union bound is not formalized in Lean.
  \item Exact logarithmic sample complexity bounds are proved in the mathematical
  explanation but represented abstractly in Lean.
  \item The project is MA-LoT-inspired but does not implement the full MA-LoT
  framework.
  \item The comparison script performs verification and structural comparison,
  not general semantic equivalence checking between arbitrary Lean files.
\end{itemize}

\section{Future Work}

Future extensions could:
\begin{itemize}[leftmargin=*]
  \item replace assumed concentration facts with Mathlib probability theorems,
  \item use real-valued risks,
  \item formalize Hoeffding's inequality,
  \item formalize exact sample complexity algebra in Lean,
  \item automate more of the MA-LoT-style repair loop,
  \item integrate an MA-LoT-related model such as a LoT-Solver model as a proof
  generation or correction agent,
  \item evaluate more theorem prompts with the same comparison pipeline.
\end{itemize}

\section{Conclusion}

This project formalizes the deterministic core of finite-class learning
guarantees in Lean 4 and makes the probability boundary explicit. The
deterministic agnostic ERM theorem is fully verified. The realizable and
agnostic high-probability results are represented using clearly named assumed
concentration results.

The project also contributes an MA-LoT-inspired proof repair study. Codex and
Claude produce proof attempts, Lean verifies them, errors are logged, repairs
are made, attempts are compared, and the final verified proof is explained for a
human reader.

The final result is both a Lean formalization and a small study of
AI-assisted theorem proving, proof repair, and proof explainability.

\appendix

\section{Appendix A: File Index}

\subsection{Core Lean Formalization}

\begin{itemize}[leftmargin=*]
  \item \texttt{LearningTheoryProject/BasicDefinitions.lean}
  \item \texttt{LearningTheoryProject/AssumedResults.lean}
  \item \texttt{LearningTheoryProject/Realizable.lean}
  \item \texttt{LearningTheoryProject/Agnostic.lean}
  \item \texttt{LearningTheoryProject/DemoExamples.lean}
  \item \texttt{LearningTheoryProject/MainTheorem.lean}
\end{itemize}

\subsection{Mathematical Notes}

\begin{itemize}[leftmargin=*]
  \item \texttt{notes/math\_overview.md}
  \item \texttt{notes/proof\_realizable.md}
  \item \texttt{notes/proof\_agnostic.md}
  \item \texttt{notes/lean\_translation\_notes.md}
\end{itemize}

\subsection{MA-LoT-Inspired Workflow}

\begin{itemize}[leftmargin=*]
  \item \texttt{malot\_workflow/README.md}
  \item \texttt{malot\_workflow/malot\_overview.md}
  \item \texttt{malot\_workflow/proof\_attempt\_log.md}
  \item \texttt{malot\_workflow/verifier\_feedback\_log.md}
  \item \texttt{malot\_workflow/repair\_iterations.md}
  \item \texttt{malot\_workflow/success\_failure\_table.md}
  \item \texttt{malot\_workflow/pattern\_analysis.md}
  \item \texttt{malot\_workflow/agnostic\_repair\_case\_study.md}
  \item \texttt{malot\_workflow/final\_reflection.md}
\end{itemize}

\subsection{Multi-Model Workflow}

\begin{itemize}[leftmargin=*]
  \item \texttt{theorem\_prompts/agnostic\_erm\_theorem.md}
  \item \texttt{theorem\_prompts/realizable\_theorem.md}
  \item \texttt{theorem\_prompts/simplification\_prompt.md}
  \item \texttt{generated\_proofs/codex/}
  \item \texttt{generated\_proofs/claude/}
  \item \texttt{multi\_model\_workflow/run\_comparison.py}
  \item \texttt{multi\_model\_workflow/comparison\_report.md}
  \item \texttt{multi\_model\_workflow/repair\_recommendations\_for\_models.md}
  \item \texttt{multi\_model\_workflow/codex\_vs\_claude\_comparison.md}
  \item \texttt{multi\_model\_workflow/consensus\_pipeline.md}
  \item \texttt{multi\_model\_workflow/reconciliation\_notes.md}
  \item \texttt{multi\_model\_workflow/semantic\_difference\_cases.md}
\end{itemize}

\subsection{Proof Refactoring}

\begin{itemize}[leftmargin=*]
  \item \texttt{proof\_refactoring/bloated\_vs\_simplified.md}
  \item \texttt{proof\_refactoring/verified\_refactors.md}
  \item \texttt{proof\_refactoring/failed\_simplifications.md}
\end{itemize}

\subsection{Presentation and Explainability}

\begin{itemize}[leftmargin=*]
  \item \texttt{presentation/final\_slides.pptx}
  \item \texttt{presentation/speaker\_notes.md}
  \item \texttt{presentation/annotated\_proof\_walkthrough.md}
  \item \texttt{presentation/paper\_vs\_lean\_table.md}
  \item \texttt{presentation/proof\_dependency\_graph.md}
  \item \texttt{presentation/proof\_pipeline\_diagram.md}
  \item \texttt{presentation/visual\_theorem\_flow.md}
  \item \texttt{presentation/final\_demo\_examples.md}
  \item \texttt{presentation/codex\_vs\_claude\_table.md}
\end{itemize}

\subsection{Final Submission Files}

\begin{itemize}[leftmargin=*]
  \item \texttt{README.md}
  \item \texttt{final\_report/final\_writeup.tex}
  \item \texttt{final\_report/final\_writeup.md}
  \item \texttt{final\_report/submission\_checklist.md}
\end{itemize}

\section{Appendix B: Main Lean Theorem Names}

\begin{itemize}[leftmargin=*]
  \item \texttt{uniform\_true\_le\_empirical\_plus}
  \item \texttt{uniform\_empirical\_le\_true\_plus}
  \item \texttt{realizable\_learning\_guarantee}
  \item \texttt{agnostic\_erm\_deterministic}
  \item \texttt{agnostic\_erm\_against\_optimal}
  \item \texttt{agnostic\_high\_probability\_guarantee}
\end{itemize}

\section{Appendix C: Final Submission Checklist}

Before submission:

\begin{enumerate}[leftmargin=*]
  \item Run \texttt{lake build}.
  \item Run \texttt{python multi\_model\_workflow/run\_comparison.py}.
  \item Confirm \texttt{presentation/final\_slides.pptx} is present.
  \item Confirm \texttt{final\_report/final\_writeup.tex} is present.
  \item Compile \texttt{final\_writeup.tex} to PDF if the submission requires PDF.
  \item Confirm no PowerPoint lock files are included. These usually begin with a tilde-dollar prefix.
  \item Commit or upload all final artifacts.
\end{enumerate}

\end{document}
